\clearpage
\SourceLeaf{404}
\begin{center}\large [GENERALISIRUNG DES LEGENDRESCHEN THEOREMS.]\end{center}
\begin{center}[I.]\end{center}
\begin{center}\textsc{Gauss} an \textsc{Olbers}. Göttingen, 9. October 1825.\end{center}
\indent\mbox{. . . . . . . . . .} Ich habe dieser Tage angefangen, in Beziehung auf mein künftiges Werk über Höhere Geodäsie, einen (sehr) kleinen Theil dessen, was die krummen Flächen betrifft, in Gedanken etwas zu ordnen. Allein ich überzeuge mich, dass ich bei der Eigenthümlichkeit meiner ganzen Behandlung des Zusammenhanges wegen gezwungen bin, \textls{sehr weit auszuholen}, so dass ich sogar meine Ansicht über die Krümmungshalbmesser bei \textls{planen} Curven vorausschicken muss. Ich bin darüber fast zweifelhaft geworden, ob es nicht gerathener sein wird, einen Theil dieser Lehren, der ganz rein geometrisch (in analytischer Form) ist und Neues mit Bekanntem gemischt in neuer Form enthält, erst besonders auszuarbeiten, es vielleicht von dem Werke abzutrennen und als eine oder zwei Abhandlungen in unsere Commentationen einzurücken. Indessen kann ich noch vorerst die Form der Bekanntmachung auf sich beruhen lassen, und werde einstweilen in dem zu Papier bringen fortfahren. \mbox{. . .}
